\documentclass[10pt,twocolumn,letterpaper]{article}
\usepackage[rebuttal]{cvpr}

% Include other packages here, before hyperref.
\usepackage{graphicx}
\usepackage{amsmath}
\usepackage{amssymb}
\usepackage{booktabs}


% If you comment hyperref and then uncomment it, you should delete
% egpaper.aux before re-running latex.  (Or just hit 'q' on the first latex
% run, let it finish, and you should be clear).
\usepackage[pagebackref,breaklinks,colorlinks,bookmarks=false]{hyperref}

% Support for easy cross-referencing
\usepackage[capitalize]{cleveref}
\crefname{section}{Sec.}{Secs.}
\Crefname{section}{Section}{Sections}
\Crefname{table}{Table}{Tables}
\crefname{table}{Tab.}{Tabs.}

% If you wish to avoid re-using figure, table, and equation numbers from
% the main paper, please uncomment the following and change the numbers
% appropriately.
%\setcounter{figure}{2}
%\setcounter{table}{1}
%\setcounter{equation}{2}

% If you wish to avoid re-using reference numbers from the main paper,
% please uncomment the following and change the counter for `enumiv' to
% the number of references you have in the main paper (here, 6).
%\let\oldthebibliography=\thebibliography
%\let\oldendthebibliography=\endthebibliography
%\renewenvironment{thebibliography}[1]{%
%     \oldthebibliography{#1}%
%     \setcounter{enumiv}{6}%
%}{\oldendthebibliography}


%%%%%%%%% COMP 472 - Project  - PLEASE UPDATE
\def\cvprPaperID{COMP 472 - Project} % *** Enter the CVPR Paper ID here
\def\confName{Project Proposal}
\def\confYear{2026}

\begin{document}

%%%%%%%%% COMP 472 - Project Proposal- PLEASE UPDATE
\title{\LaTeX\ COMP 472 - Project Proposal}  % **** Enter the paper title here


\maketitle
\thispagestyle{empty}
\appendix

%%%%%%%%% BODY TEXT - ENTER YOUR RESPONSE BELOW
\section{Team}

\begin{enumerate}
	\item[] Lucas Graham 40249532
	\item[] Dominique Proulx 40177566
	\item[] Valeriia Nikandrova 40157880
	\item[] Azmi Abidi 40248132
	\item[] Kalin Milanov 40179417
\end{enumerate}


%-------------------------------------------------------------------------

\section{Problem Statement and Application}
\textcolor{red}{Provide a background about the topic and specify why the problem is important?
What are the associated challenges of your selected problem application? What are your expectations/goals
throughout developing the application of interest?}
%------------------------------------------------------------------------
\section{Image Dataset Selection}


\textcolor{red}{Explain your image dataset(s) and provide the statistical details of your data (e.g. number of
images, image size, image format, etc). Specify where you have found the dataset(s) and means of access to the data
(e.g. published paper, downlink, etc). You may consult Kaggle to find possible image datasets of interest, or you can
use any data you have collected from your previous research activity or have access to.}

\subsection{Dataset 1}
\subsection{Dataset 2}
\subsection{Dataset 3}

%-------------------------------------------------------------------------
\section{Possible Methodology}

{\textcolor{red}{Highlight the “possible methods” you could use to solve the problem. Specify how you will be
handling/processing the data to train your deep learning pipeline using a CNN model. Furthermore, discuss the
metrics that will be used to assess and evaluate the pipeline, and your expectations regarding the kind of
results/performance to be achieved. You need to discuss the possible method(s) and how the obtained results will be
compared and analyzed to each other. Further, you need to discuss the potentials of your analysis and comparisons
and how they can be useful for scientists and engineering in the field. }
%------------------------------------------------------------------------
\section{Gantt Chart}
{\textcolor{red}{On 2nd page
use an additional page (supplemental material) to illustrate a Gantt chart of the project development to
list (a) schedules and (b) items of milestones and deliverables. Note that you cannot use this page to extend your
proposal description. }

{\textcolor{red}When placing figures in \LaTeX, it is almost always best to use \verb+\includegraphics+, and to specify the  figure width as a multiple of the line width as in the example below
{\small\begin{verbatim}
   \usepackage{graphicx} ...
   \includegraphics[width=0.8\linewidth]
                   {myfile.pdf}
\end{verbatim}
}
\begin{figure}[t]
	\centering
	\fbox{\rule{0pt}{0.5in} \rule{0.9\linewidth}{0pt}}
	%\includegraphics[width=0.8\linewidth]{egfigure.eps}
	\caption{Example of caption.  It is set in Roman so that mathematics
		(always set in Roman: $B \sin A = A \sin B$) may be included without an
		ugly clash.}
	\label{fig:onecol}
\end{figure}

%%%%%%%%% REFERENCES
{\textcolor{red}{You can add an additional page (if needed) to extend your reference list cited in your proposal. The
citations may include, but not limited to, published papers and domain links. (include a link to your dataset). Please
note that failure to properly cite your references constitutes to a plagiarism and will be deemed for reporting}
{\small
\bibliographystyle{ieee_fullname}
\bibliography{egbib}
}

\end{document}
